\documentclass[12pt]{article}
\usepackage[T1]{fontenc}
\usepackage[utf8]{inputenc}
\usepackage{lmodern}
\usepackage{textcomp}
\usepackage{enumitem}
\usepackage{geometry}
\geometry{margin=1in}
\begin{document}
\begin{center}
Answer Key - Astronomy - Graduate Level Assessment 15 \
\small (Answers are ordered exactly as in the question paper.)
\end{center}

\section*{Multiple Choice}
\begin{enumerate}[label=\arabic*.]
\item \textbf{Answer:} C
\\ \textit{Options:} A) about 18 years; B) 24 hours; C) 29 Earth days; D) a year
\\ \textit{Note:} Correct option: C - 29 Earth days
\item \textbf{Answer:} C
\\ \textit{Options:} A) The Moon will be closer to the Earth and the Earth's day will be longer.; B) The Moon will be closer to the Earth and the Earth's day will be shorter.; C) The Moon will be further from the Earth and the Earth's day will be longer.; D) The Moon will be further from the Earth and the Earth's day will be shorter.
\\ \textit{Note:} Correct option: C - The Moon will be further from the Earth and the Earth's day will be longer.
\item \textbf{Answer:} A
\\ \textit{Options:} A) 750 million years.; B) 2 billion years.; C) 3 billion years.; D) 4.5 billion years.
\\ \textit{Note:} Correct option: A - 750 million years.
\item \textbf{Answer:} D
\\ \textit{Options:} A) Titan is the only outer solar system moon with a thick atmosphere; B) Titan is the only outer solar system moon with evidence for recent geologic activity; C) Titan's atmosphere is composed mostly of hydrocarbons; D) A and D
\\ \textit{Note:} Correct option: D - A and D
\item \textbf{Answer:} C
\\ \textit{Options:} A) The first time you drop a bowling ball and it falls to the ground proving your hypothesis.; B) After you've repeated your experiment many times.; C) You can never prove your theory to be correct only "yet to be proven wrong".; D) When you and many others have tested the hypothesis.
\\ \textit{Note:} Correct option: C - You can never prove your theory to be correct only "yet to be proven wrong".
\end{enumerate}

\section*{True / False}
\begin{enumerate}[label=\arabic*.]
\setcounter{enumi}{5}
\item \textbf{Answer:} True
\\ \textit{Options:} A) True; B) False
\\ \textit{Note:} LLM-generated true statement based on MCQ.
\item \textbf{Answer:} False
\\ \textit{Options:} A) True; B) False
\\ \textit{Note:} LLM-generated false statement. Correct answer: 'These craters contain the only permanently shadowed regions on Mercury'.
\item \textbf{Answer:} True
\\ \textit{Options:} A) True; B) False
\\ \textit{Note:} LLM-generated true statement based on MCQ.
\item \textbf{Answer:} False
\\ \textit{Options:} A) True; B) False
\\ \textit{Note:} LLM-generated false statement. Correct answer: 'black holes.'.
\item \textbf{Answer:} True
\\ \textit{Options:} A) True; B) False
\\ \textit{Note:} LLM-generated true statement based on MCQ.
\end{enumerate}

\section*{Long Form Response}
\begin{enumerate}[label=\arabic*.]
\setcounter{enumi}{10}
\item \textbf{Answer:} Solar eclipses would be much more frequent.. Explain why this is the correct answer and provide supporting evidence. Reference key facts or concepts that support this choice.
\\ \textit{Note:} Long-form derived from MCQ; award credit for stating the correct idea and giving supporting reasoning.
\item \textbf{Answer:} 300 Kelvin. Explain why this is the correct answer and provide supporting evidence. Reference key facts or concepts that support this choice.
\\ \textit{Note:} Long-form derived from MCQ; award credit for stating the correct idea and giving supporting reasoning.
\end{enumerate}

\vfill
\noindent\textit{End of Answer Key}
\end{document}